\subsection{20.100: distinct powers, reduction, and four finite proofs}
\label{partial:20.100}
\problemstatement{20.100}

For a positive integer \(n\), assume \(G\) has no element of order
\(2,\ldots,n+1\). The question asks whether every \(n\)-element set
can be ordered as \(a_1,\ldots,a_n\) with
\(a_1,a_2^2,\ldots,a_n^n\) pairwise distinct. The identity is allowed.
Sun's Conjecture~4.1 is this assertion; his Theorem~1.4 covers
torsion-free abelian groups and \(n\le3\) \cite{sun-permutations}.
The experiment obtained a torsion-free theorem without commutativity,
a reduction to finite abelian groups, an explicit finite
counterexample bound for fixed \(n\), and complete certificates
for \(n=4,5,6,7\).

\begin{proposition}
Any finite list of infinite-order elements, with repetitions allowed,
can be assigned prescribed distinct positive integer exponents so
that the resulting powers are distinct.
\end{proposition}
\begin{proof}
On the cyclic subgroups generated by the list, nontrivial
intersection is an equivalence relation: intersect the two
finite-index subgroups of the middle infinite cyclic group to
prove transitivity. In each of the finitely many classes choose
a nonidentity common power \(c\). Write \(c=a^{k_a}\), \(k_a\ne0\),
for each entry in that class and give it length \(1/|k_a|\).
If \(a^i=b^j\), then the entries lie in the same class, and raising
to the integer \(k_ak_b\) gives \(c^{ik_b}=c^{jk_a}\).
Infinite order implies \(i/|k_a|=j/|k_b|\).
Sort the entire list by nondecreasing positive length and assign
the exponents in increasing order. The products of exponent and
length strictly increase, so no equality is possible.
One additional identity element can be assigned any unused exponent.
\end{proof}

For fixed \(n\), the question for all groups is equivalent to the
question for finite abelian groups. To see the nontrivial direction,
partition the finite-order nonidentity entries according to the
cyclic subgroups they generate. All prime divisors of their orders
exceed \(n+1\); otherwise \(G\) has an element of forbidden prime
order. Thus \(\langle a^i\rangle=\langle a\rangle\) for
\(1\le i\le n\). Distinct cyclic groups cannot contribute equal
powers in this exponent range.
Represent the entries in each cyclic group in a separate abstract
cyclic coordinate. This preserves exactly their relevant power
equalities. Replace all infinite-order entries by distinct
nonidentity elements in one additional cyclic coordinate of prime
order greater than \(n+1\), and represent the identity by the identity.
An affirmative finite abelian result assigns the exponents.
Keep the assignments on finite-order entries; reassign the
remaining exponents to the infinite-order entries using the
proposition. Finite- and infinite-order powers cannot coincide.
This proves the equivalence, without embedding \(G\) in an abelian
group.

There is also a finite size bound. Let \(p_n\) be the least prime
greater than \(n+1\). A counterexample at \(n\ge2\), if one exists,
exists in a finite abelian group of order at most
\[
 p_n n^{n-1}.
\]
Indeed, from a finite abelian counterexample choose a collision
for every exponent permutation and impose those relations
\(ix_a=jx_b\) on a free module over
\(S=\Z[1/P]\), where \(P\) is the product of primes at most \(n+1\).
Its distinguished elements remain distinct by mapping to the
counterexample, and every permutation fails. In the multigraph
of relations, a spanning tree in a component of \(k\) vertices
expresses each variable as an \(S\)-unit multiple of one generator.
A nontree edge either gives the zero relation or gives an
annihilator, up to an \(S\)-unit, equal to the difference of the
two coefficient products around a simple cycle. A nonzero difference
has absolute value less than \(n^k\).
The component is therefore free of rank one over \(S\), or cyclic
finite of order less than \(n^k\), after stripping inverted primes.
If every component is finite, the total order is at most \(n^n<
p_nn^{n-1}\). Otherwise the finite components involve at most
\(n-1\) variables; replace all infinite-order variables by a separate
\(C_{p_n}\) as above. A good assignment would lift back to a
contradiction, and the resulting order is at most \(p_nn^{n-1}\).
This is decidability for each fixed \(n\), not a uniform proof in \(n\).

\paragraph{The certificate is a proof of an all-group assertion.}
A node consists of an integer relation matrix \(R\) on \(n\)
distinguished variables, interpreted over \(S\). It asserts that
no counterexample tuple satisfies those relations. There are two
rules. A leaf proves that a coordinate difference belongs to the
\(S\)-row span of \(R\), forcing two entries equal. A branch chooses
one exponent permutation and treats every one of its
\(\binom n2\) possible collisions. For each collision it points
to an earlier proved node, with a bijective variable relabeling,
and proves that every relation of that child follows from the
parent relations and the collision. The children need not describe
exact quotients; the indicated implications suffice.
The final node has no relations. Induction along earlier-node
references proves that no counterexample exists in any \(S\)-module,
hence in any eligible finite abelian group and, by the reduction,
in any eligible group.

The producer represents connected components by rational weights
and annihilators and uses gcd operations to merge collisions.
The GAP checker instead uses integer Smith normal form:
it verifies \(UMV=D\), unimodularity of \(U,V\), and diagonality
of \(D\), then tests divisibility after removing exactly the
inverted prime factors and requires zero in all free coordinates.
This checks the mathematical implications without trusting the
producer's normalization. A separate data-only parser checks the
certificate grammar, excluding executable instructions that could
alter the checker.

\begin{center}\small
\begin{tabular}{@{}rrrr@{}}
\toprule
\(n\)&Nodes&Collision cases&Relation implications\\\midrule
4&227&1,056&3,953\\
5&10,493&87,570&426,345\\
6&265,104&3,245,190&19,094,383\\
7&42,891,332&786,577,344&5,447,182,055\\
\bottomrule
\end{tabular}
\end{center}

For \(n=7\), six GAP workers partitioned identifiers by residue
modulo six. Each read all literal states but checked its assigned
nodes; together their disjoint classes cover the proof, with
every child identifier smaller than its parent. All six recorded
actual process exit zero. The completed aggregate was written
at 13:08:05 UTC on 12 September, about 26 hours after launch.
The complete compressed certificate occupies 1,271,256,410 bytes.
The six 14\,GiB workspace ceilings are configuration limits,
not a measurement of peak memory.
These historical completion records have been inspected in the
manuscript phase; the 26-hour verification has not been represented
as a fresh paper-phase replay.

As a hypothesis control, allowing order five at \(n=4\) yields
the four nonidentity elements of \(C_5\). In additive coordinates
modulo five, distinct nonzero outputs would have product \(-1\),
whereas the input product times the product of exponents is \(1\).
This is a counterexample to the weakened assumption and is stored
as such. Sun's higher-\(n\) small cyclic searches are a separate
kind of evidence from these universal fixed-\(n\) certificates.
The unrestricted conjecture remains open in this investigation.
Detailed proofs, grammar and receipts are indexed by
\artifact{research/20.100-reduction.md} and
\artifact{research/20.100-n7-report.md}.
